\documentclass[11pt,a4paper]{article}
\usepackage[utf8]{inputenc}
\usepackage[T1]{fontenc}
\usepackage{geometry}
\usepackage{hyperref}
\usepackage{graphicx}
\usepackage{listings}
\usepackage{xcolor}
\usepackage{enumitem}
\usepackage{titlesec}
\usepackage{fancyhdr}

% Configuration
\geometry{margin=1in}
\hypersetup{
    colorlinks=true,
    linkcolor=blue,
    filecolor=magenta,      
    urlcolor=cyan,
    pdftitle={Bastion Password Manager - User Guide},
}

% Header and Footer
\pagestyle{fancy}
\fancyhf{}
\rhead{Bastion Password Manager}
\lhead{User Guide}
\rfoot{Page \thepage}

% Title Page
\title{\textbf{Bastion Password Manager} \\ User Guide}
\author{Bastion Development Team}
\date{Version 1.0 | November 2025}

\begin{document}

\maketitle
\thispagestyle{empty}
\tableofcontents
\newpage

\section{Introduction}
Welcome to Bastion! This guide will help you get started with your secure password manager.

\section{Installation}

\subsection{Desktop Application}

\subsubsection{Windows}
\begin{enumerate}
    \item Download \texttt{Bastion Setup.exe}.
    \item Double-click the installer.
    \item Follow the installation wizard.
    \item Launch Bastion from the Start menu or desktop shortcut.
\end{enumerate}

\subsubsection{macOS}
\begin{enumerate}
    \item Download \texttt{Bastion.dmg}.
    \item Open the DMG file.
    \item Drag Bastion to your Applications folder.
    \item Open Bastion from Applications.
\end{enumerate}

\subsubsection{Linux}
\begin{enumerate}
    \item Download \texttt{Bastion.AppImage}.
    \item Make it executable: \texttt{chmod +x Bastion.AppImage}.
    \item Run it: \texttt{./Bastion.AppImage}.
\end{enumerate}


\section{Getting Started}

\subsection{Creating Your Account}
\begin{enumerate}
    \item \textbf{Launch Bastion} desktop application.
    \item Click \textbf{"Sign Up"}.
    \item Enter your \textbf{email address}.
    \item Create a \textbf{strong master password}:
    \begin{itemize}
        \item Use at least 12 characters.
        \item Mix uppercase, lowercase, numbers, and symbols.
        \item Make it memorable but unique.
        \item \textbf{IMPORTANT}: This password cannot be recovered if forgotten!
    \end{itemize}
    \item Click \textbf{"Create Account"}.
\end{enumerate}
You're now ready to start using Bastion!

\subsection{Your First Login}
\begin{enumerate}
    \item \textbf{Launch Bastion}.
    \item Enter your \textbf{email} and \textbf{master password}.
    \item Click \textbf{"Log In"}.
    \item Your vault will unlock.
\end{enumerate}
\textbf{Tip}: The desktop app will keep you logged in for 24 hours by default.

\section{Managing Your Vault}

\subsection{Adding a Password}
\begin{enumerate}
    \item In the Bastion app, go to the \textbf{"Vault"} section.
    \item Click \textbf{"Add Entry"}.
    \item Fill in the details:
    \begin{itemize}
        \item \textbf{Website URL}: \texttt{https://example.com}
        \item \textbf{Username}: Your username or email
        \item \textbf{Password}: Your password (or generate a new one)
        \item \textbf{Notes} (optional): Additional information
    \end{itemize}
    \item Click \textbf{"Save"}.
\end{enumerate}
Your password is now encrypted and stored securely!

\subsection{Editing a Password}
\begin{enumerate}
    \item Find the entry in your vault.
    \item Click on it to view details.
    \item Click \textbf{"Edit"}.
    \item Make your changes.
    \item Click \textbf{"Save"}.
\end{enumerate}

\subsection{Deleting a Password}
\begin{enumerate}
    \item Find the entry in your vault.
    \item Click on it to view details.
    \item Click \textbf{"Delete"}.
    \item Confirm deletion.
\end{enumerate}
\textbf{Warning}: Deleted passwords cannot be recovered.

\subsection{Searching Your Vault}
Use the search bar at the top to quickly find entries by:
\begin{itemize}
    \item Website name
    \item URL
    \item Username
\end{itemize}


\section{Password Generator}
Bastion includes a smart password generator that creates strong, site-compliant passwords.

\subsection{Using the Generator}

\subsubsection{Method 1: In the Desktop App}
\begin{enumerate}
    \item When adding or editing an entry, click \textbf{"Generate Password"}.
    \item Review the generated password.
    \item Adjust settings if needed:
    \begin{itemize}
        \item Length (12-64 characters)
        \item Include uppercase letters
        \item Include lowercase letters
        \item Include numbers
        \item Include symbols
    \end{itemize}
    \item Click \textbf{"Use This Password"}.
\end{enumerate}


\subsection{Password Strength Indicator}
\begin{itemize}
    \item \textbf{Weak} (Red): Too short or simple
    \item \textbf{Fair} (Orange): Decent but could be better
    \item \textbf{Good} (Yellow): Strong enough
    \item \textbf{Excellent} (Green): Very strong
\end{itemize}
Always aim for "Good" or "Excellent" strength.

\section{Family Sharing}
Share passwords with family members securely—they can use the passwords without seeing them!

\subsection{Understanding Family Sharing}
\begin{itemize}
    \item \textbf{Owner}: You create the family and share credentials.
    \item \textbf{Members}: Family members can autofill shared passwords but \textbf{cannot view them}.
    \item \textbf{Security}: Shared passwords are never visible to members—only autofilled.
\end{itemize}

\subsection{Creating a Family}
\begin{enumerate}
    \item In Bastion, go to \textbf{"Families"} section.
    \item Click \textbf{"Create Family"}.
    \item Enter a family name (e.g., "Smith Family").
    \item Click \textbf{"Create"}.
\end{enumerate}
You are now the owner of this family!

\subsection{Inviting Family Members}
\begin{enumerate}
    \item Open your family.
    \item Click \textbf{"Create Invite Link"}.
    \item Set expiration (default: 7 days).
    \item \textbf{Copy the link} and send it to your family member via email, messaging app, etc.
    \item They click the link and accept the invite.
\end{enumerate}
\textbf{Note}: Invite links expire and can only be used once.

\subsection{Sharing a Password}
\begin{enumerate}
    \item Go to your \textbf{"Vault"}.
    \item Select the credential you want to share.
    \item Click \textbf{"Share with Family"}.
    \item Choose which family to share with.
    \item \textbf{(Optional)} Set allowed domains:
    \begin{itemize}
        \item Leave empty to allow all domains.
        \item Or specify: \texttt{netflix.com}, \texttt{facebook.com}.
    \end{itemize}
    \item Click \textbf{"Share"}.
\end{enumerate}
The password is now shared! Family members can autofill it.

\subsection{Using Shared Passwords (As a Member)}
\begin{enumerate}
    \item Log in to the desktop app.
    \item Navigate to the login page (e.g., netflix.com).
    \item Copy the username and password from Bastion.
    \item Paste into the login form.
\end{enumerate}

\textbf{What You'll See:}
\begin{itemize}
    \item Website name and URL
    \item Username (visible)
    \item Password: \texttt{••••••••••} (masked)
\end{itemize}

\textbf{What You Can Do:}
\begin{itemize}
    \item Autofill the password
    \item Copy the username and password (using the copy buttons)
\end{itemize}

\textbf{What You Cannot Do:}
\begin{itemize}
    \item View the actual password (it remains masked)
    \item Edit the password
    \item Share it with others
\end{itemize}

\subsection{Viewing Who Used Your Shared Passwords}
\begin{enumerate}
    \item Go to \textbf{"Families"} $\rightarrow$ Select your family.
    \item Click \textbf{"Audit Log"}.
    \item View recent activity:
    \begin{itemize}
        \item Who requested autofill
        \item Which password
        \item Which website
        \item When (timestamp)
        \item Success or failure
    \end{itemize}
\end{enumerate}

\subsection{Stopping Sharing}
\begin{enumerate}
    \item Go to \textbf{"Families"} $\rightarrow$ Select your family.
    \item Click \textbf{"Shared Credentials"}.
    \item Find the credential.
    \item Click \textbf{"Unshare"} or \textbf{"Revoke"}.
    \item Confirm.
\end{enumerate}
Members will immediately lose access to this password.

\subsection{Removing a Family Member}
\begin{enumerate}
    \item Go to \textbf{"Families"} $\rightarrow$ Select your family.
    \item Click \textbf{"Members"}.
    \item Select the member to remove.
    \item Click \textbf{"Remove from Family"}.
    \item Confirm.
\end{enumerate}
They will lose access to all shared credentials immediately.

\section{Account Recovery}
If you forget your master password, you can recover your account through trusted contacts.

\subsection{Setting Up Account Recovery}

\subsubsection{Step 1: Add Trusted Contacts}
\begin{enumerate}
    \item Go to \textbf{"Account Recovery"} section.
    \item Click \textbf{"Add Trusted Contact"}.
    \item Choose a contact from:
    \begin{itemize}
        \item Your family members (if in a family)
        \item Other Bastion users (by email)
    \end{itemize}
    \item Click \textbf{"Add"}.
\end{enumerate}
\textbf{Recommendation}: Add 2-3 trusted contacts for redundancy.

\subsubsection{Step 2: They Accept}
\begin{enumerate}
    \item Your contact receives a notification or request.
    \item They review and accept being your trusted contact.
    \item System encrypts your master key for them.
\end{enumerate}
\textbf{Note}: Trusted contacts can help you recover but \textbf{cannot access your vault} unless you request recovery.

\subsection{Recovering Your Account}
If you forget your master password:

\subsubsection{Step 1: Initiate Recovery}
\begin{enumerate}
    \item On the login screen, click \textbf{"Forgot Password?"}.
    \item Enter your \textbf{email address}.
    \item Select a \textbf{trusted contact} from the list.
    \item Click \textbf{"Request Recovery"}.
\end{enumerate}

\subsubsection{Step 2: Wait for Approval}
\begin{enumerate}
    \item Your trusted contact will receive a notification.
    \item They log into Bastion.
    \item They go to \textbf{"Account Recovery"} section.
    \item They see your request and click \textbf{"Approve"}.
    \item They enter their master password to confirm.
\end{enumerate}

\subsubsection{Step 3: Complete Recovery}
\begin{enumerate}
    \item Check the status in your recovery request.
    \item Once approved, click \textbf{"Complete Recovery"}.
    \item Enter a \textbf{new master password}.
    \item Click \textbf{"Reset Password"}.
\end{enumerate}
Your vault is now accessible with your new password!

\textbf{Important Notes:}
\begin{itemize}
    \item The trusted contact \textbf{cannot access your vault}.
    \item They only approve the recovery.
    \item You must set a \textbf{new} master password.
    \item Your vault data remains encrypted throughout.
\end{itemize}

\section{Security Best Practices}

\subsection{Master Password}
\textbf{DO:}
\begin{itemize}
    \item Use a unique password (not used elsewhere).
    \item Make it at least 12 characters.
    \item Use a passphrase: "correct-horse-battery-staple".
    \item Write it down and store it securely (physical safe).
\end{itemize}

\textbf{DON'T:}
\begin{itemize}
    \item Use common passwords.
    \item Reuse passwords from other accounts.
    \item Share your master password with anyone.
    \item Store it digitally (in emails, notes, etc.).
\end{itemize}

\subsection{Vault Passwords}
\textbf{DO:}
\begin{itemize}
    \item Use the password generator for new passwords.
    \item Use unique passwords for each site.
    \item Use the maximum allowed length.
    \item Enable all character types (uppercase, lowercase, numbers, symbols).
    \item Update old/weak passwords regularly.
\end{itemize}

\textbf{DON'T:}
\begin{itemize}
    \item Reuse passwords across sites.
    \item Use personal information (birthdays, names).
    \item Use simple patterns (123456, qwerty).
\end{itemize}

\subsection{Family Sharing}
\textbf{DO:}
\begin{itemize}
    \item Only share with people you completely trust.
    \item Review audit logs regularly.
    \item Use domain restrictions when possible.
    \item Revoke access when someone leaves.
    \item Set up multiple trusted contacts for recovery.
\end{itemize}

\textbf{DON'T:}
\begin{itemize}
    \item Share your master password.
    \item Share credentials you don't want accessed.
    \item Ignore unusual activity in audit logs.
    \item Trust screenshots of "shared" passwords.
\end{itemize}

\subsection{General Security}
\textbf{DO:}
\begin{itemize}
    \item Lock Bastion when stepping away.
    \item Log out on shared computers.
    \item Keep the desktop app updated.
    \item Set up account recovery before you need it.
\end{itemize}

\textbf{DON'T:}
\begin{itemize}
    \item Leave Bastion unlocked and unattended.
    \item Use Bastion on untrusted/public computers.
    \item Ignore security warnings or prompts.
    \item Share invite links publicly.
\end{itemize}

\section{Troubleshooting}

\subsection{I Forgot My Master Password}
See Account Recovery section. If you didn't set up trusted contacts, \textbf{your account cannot be recovered}. This is by design for security.


\subsection{"Cannot Connect to Server" Error}
\textbf{Possible Causes:}
\begin{itemize}
    \item Backend server is down.
    \item Network connection issue.
    \item Firewall blocking connection.
\end{itemize}

\textbf{Solution:}
\begin{enumerate}
    \item Check your internet connection.
    \item Contact your system administrator.
    \item Wait a few minutes and try again.
\end{enumerate}

\subsection{Shared Password Not Showing Up}
\textbf{For Members:}
\begin{enumerate}
    \item Make sure you accepted the family invite.
    \item Refresh your vault (log out and log back in).
    \item Check with the owner that they shared it.
\end{enumerate}

\textbf{For Owners:}
\begin{enumerate}
    \item Verify you clicked "Share".
    \item Check the family has the right members.
    \item Check audit logs for any errors.
\end{enumerate}

\subsection{Desktop App Won't Start}
\textbf{Windows:}
\begin{enumerate}
    \item Try running as Administrator.
    \item Check Windows Defender didn't block it.
    \item Reinstall the application.
\end{enumerate}

\textbf{macOS:}
\begin{enumerate}
    \item Go to System Preferences $\rightarrow$ Security \& Privacy.
    \item Allow Bastion to run.
    \item If still failing, reinstall.
\end{enumerate}

\textbf{Linux:}
\begin{enumerate}
    \item Make sure the AppImage is executable: \texttt{chmod +x Bastion.AppImage}.
    \item Install FUSE if needed: \texttt{sudo apt install fuse}.
    \item Try extracting and running directly.
\end{enumerate}

\section{FAQs}

\subsection{General Questions}
\textbf{Q: Is my data safe?} \\
A: Yes! Your master key never leaves your device. All data is encrypted before being sent to the server. Even if the server is compromised, your passwords remain secure.

\textbf{Q: Can Bastion employees see my passwords?} \\
A: No. We use zero-knowledge encryption. Your passwords are encrypted on your device, and we never have access to your master key.

\textbf{Q: What happens if I lose my master password?} \\
A: If you set up trusted contacts, you can recover your account. Otherwise, there is no way to recover your account—this is by design for security.

\textbf{Q: Can I use Bastion offline?} \\
A: The desktop app requires internet connection to sync. However, once synced, you can view (but not edit) your vault offline.

\textbf{Q: How much does Bastion cost?} \\
A: Contact your administrator for pricing/licensing information.

\subsection{Family Sharing}
\textbf{Q: Can family members see my passwords?} \\
A: No! Family members can autofill or copy shared passwords, but they remain masked in the UI. They cannot view the plaintext password.

\textbf{Q: Can family members share credentials they received?} \\
A: No. Only the owner can share credentials.

\textbf{Q: How many family members can I have?} \\
A: There is no hard limit, but we recommend keeping families small and trustworthy.

\textbf{Q: Can a shared password be used on any website?} \\
A: Only if you allow it. You can restrict shared passwords to specific domains when sharing.

\textbf{Q: Will I know if a family member uses a shared password?} \\
A: Yes! Every autofill request is logged in the audit log with timestamp and website.

\subsection{Technical Questions}
\textbf{Q: What encryption does Bastion use?} \\
A: We use XChaCha20-Poly1305 for encryption and Argon2id for password hashing.

\textbf{Q: Where is my data stored?} \\
A: Encrypted data is stored on the backend server. Your master key and decrypted passwords only exist on your device.

\textbf{Q: Can I export my passwords?} \\
A: This feature is coming soon. Contact your administrator for manual export if needed.

\textbf{Q: Does Bastion work on mobile?} \\
A: Not yet. Mobile apps are planned for a future release.

\textbf{Q: Can I import passwords from another password manager?} \\
A: This feature is coming soon.


\section{Getting Help}
If you need additional assistance:
\begin{itemize}
    \item \textbf{Documentation}: Reread this guide and the README.
    \item \textbf{System Administrator}: Contact your IT department or Bastion administrator.
    \item \textbf{Email Support}: support@bastion-pm.example.com (update with actual contact).
    \item \textbf{Report Bugs}: Contact your administrator with details.
\end{itemize}

\section{Tips for Success}
\begin{enumerate}
    \item \textbf{Set up trusted contacts immediately} after creating your account.
    \item \textbf{Use the password generator} for all new accounts.
    \item \textbf{Review shared credentials} weekly if you're a family owner.
    \item \textbf{Update passwords regularly}, especially after security breaches.
    \item \textbf{Lock your vault} when stepping away from your computer.
    \item \textbf{Never share your master password} with anyone, including family.
\end{enumerate}

\end{document}
